\chapter{The Runtime Block}
\label{ch:runtime-block}

\newthought{The optional \yamlkey{Runtime} block} tunes \emph{how} a scan executes---how many
workers, how points are batched, whether per-sample folders are written, and how external commands
are scheduled. You never need it to get a correct result; you reach for it to go faster or to
control disk use on large runs.

\section{The keys}
\label{sec:rt-keys}

\yamlkey{Runtime} is a top-level block (Table~\ref{tab:rt-keys}). Every key has a sensible default,
so an empty or absent block behaves exactly as before.

\begin{table}[ht]
  \footnotesize
  \begin{tabular}{@{}p{0.33\linewidth}lp{0.40\linewidth}@{}}
    \toprule
    \textbf{Key} & \textbf{Default} & \textbf{Meaning} \\
    \midrule
    \yamlkey{mode}             & \code{auto}  & execution mode: \code{auto}, \code{thread}, or \code{process} \\
    \yamlkey{writer}           & \code{auto}  & \textsc{hdf5} writer: \code{auto}, \code{thread}, or \code{process} \\
    \yamlkey{workers}          & \code{0}     & process-worker count (\code{0} = choose automatically) \\
    \yamlkey{batch\_size}      & \code{256}   & points proposed and persisted per batch \\
    \yamlkey{writer\_record\_batch\_size}     & \code{100} & process writer: records micro-batched before a flush \\
    \yamlkey{writer\_record\_flush\_interval} & \code{5.0} & process writer: max seconds between flushes \\
    \yamlkey{sample\_artifacts} & \code{auto} & per-sample folders: \code{auto}, \code{always}, \code{never} \\
    \yamlkey{Subprocess}       & ---          & external-command scheduler knobs (below) \\
    \bottomrule
  \end{tabular}
  \caption{\yamlkey{Runtime} keys and defaults. Every key is optional, and each defaults to
  v1-equivalent behaviour.}
  \label{tab:rt-keys}
\end{table}

\begin{yamlwin}
Runtime:
  mode: auto           # thread today; opt into processes with: process
  workers: 0           # 0 = auto (min(cpu, 4))
  batch_size: 256      # rows per group; 1 reproduces v1 exactly
  writer: auto         # follow the execution mode
  sample_artifacts: auto
\end{yamlwin}

\section{Execution mode: thread and process}
\label{sec:rt-mode}

\yamlkey{mode} chooses how worker points are executed:

\begin{description}[style=nextline, leftmargin=1.2em]
  \item[\code{auto}] (default) the v1-compatible single-process \emph{thread} path. Today
        \code{auto} always resolves to \code{thread}---process workers are strictly opt-in.
  \item[\code{thread}] force the thread path.
  \item[\code{process}] opt into multiple worker \emph{processes}, which scale throughput with
        \textsc{cpu} cores. Available for \yamlkey{Operas}-only scans (see below).
\end{description}

\noindent \yamlkey{workers} sets the process count under \code{mode: process}; \code{0} (the default)
chooses \code{min(cpu\_count, 4)}, then caps it at \yamlkey{make\_paraller}.

\begin{jwarn}
\code{mode: auto} is \emph{not} process mode---auto is thread today. To run on processes you must
write \code{mode: process} explicitly.
\end{jwarn}

\Jarvis\ enables \code{mode: process} only when it is safe to. If \emph{any} of the following hold,
the scan silently and safely runs as \code{thread} instead---the reason is logged, never a crash:

\begin{itemize}
  \item the workflow contains a \yamlkey{Calculators} module (external calculators stay thread-only
        for now);
  \item any module references the \token{@Sdir} token;
  \item \yamlkey{sample\_artifacts} is \code{always};
  \item the normalized config cannot be pickled out to the workers;
  \item the platform has no \textsc{posix} shared memory.
\end{itemize}

\begin{jnote}
Process mode is reproducible. At a fixed seed the persisted records are identical whether you run
with one worker or many, \textsc{uuid}s stay unique, and \cli{-{}-resume} never duplicates points:
each worker draws a deterministic child random stream from a master seed stored in the checkpoint.
\end{jnote}

\section{Batching}
\label{sec:rt-batch}

\yamlkey{batch\_size} sets how many points \Jarvis\ proposes, executes, and persists as one
group---default \code{256}, rather than one at a time. Batching the data plane cuts per-point
overhead and keeps workers fed; larger batches amortise more overhead at the cost of coarser
checkpoint granularity. \code{batch\_size: 1} reproduces the exact v1 per-sample sequence (it is
golden-tested), occasionally useful for bit-for-bit comparisons. Tuning is the subject of
Chapter~\ref{ch:performance}.

\section{The HDF5 writer}
\label{sec:rt-writer}

\yamlkey{writer} chooses where persistence runs. \code{auto} (default) follows the execution
mode---a separate writer \emph{process} when \yamlkey{mode} resolves to process and shared memory is
available, otherwise an in-process writer \emph{thread}. Force either with \code{thread} or
\code{process}. Two keys tune the process writer, and \emph{only} the process writer:

\begin{description}[style=nextline, leftmargin=1.2em]
  \item[\yamlkey{writer\_record\_batch\_size}] (default \code{100}) how many single-record hand-offs
        to micro-batch before a flush.
  \item[\yamlkey{writer\_record\_flush\_interval}] (default \code{5.0}) the longest a record waits,
        in seconds, before it is flushed.
\end{description}

\begin{jwarn}
The \yamlkey{writer\_record\_*} keys do nothing unless the writer actually runs as a process---under
the default thread writer they are ignored.
\end{jwarn}

\section{Per-sample artifacts}
\label{sec:rt-artifacts}

\yamlkey{sample\_artifacts} decides whether \Jarvis\ creates a working folder on disk for every
sample---which, for fast in-process scans, can dominate the runtime:

\begin{description}[style=nextline, leftmargin=1.2em]
  \item[\code{auto}] (default) skip per-sample folders for opera-only scans, where they are not
        needed; create them when a workflow actually needs them.
  \item[\code{always}] always create per-sample folders (the classic behaviour).
  \item[\code{never}] never create them, even to record a failure.
\end{description}

\begin{jnote}
Using \token{@Sdir} in a path, or running any \yamlkey{Calculators} workflow, forces a sample to
\emph{materialize} its folder regardless of \code{auto}. So \code{auto} only skips disk work when
it is genuinely safe to---typically pure-\yamlkey{Operas} scans. Failed lazy samples still get a
replayed log (Chapter~\ref{ch:logging}).
\end{jnote}

\section{The subprocess scheduler}
\label{sec:rt-subprocess}

External calculator commands run through a shared subprocess scheduler whose knobs live under
\yamlkey{Runtime.Subprocess}:

\begin{yamlwin}
Runtime:
  Subprocess:
    max_concurrency: 8
    max_pending: 256
    per_task_timeout_sec: 300
    terminate_grace_sec: 5
    log_policy: logger
    progress_interval_sec: 5
    diagnostics_enabled: false
    diagnostics_interval_sec: 10
\end{yamlwin}

\begin{description}[style=nextline, leftmargin=1.2em]
  \item[\yamlkey{max\_concurrency}] how many external commands may run at once.
  \item[\yamlkey{max\_pending}] how many may queue waiting for a slot.
  \item[\yamlkey{per\_task\_timeout\_sec}] a per-command time limit.
  \item[\yamlkey{terminate\_grace\_sec}] grace period before a timed-out command is killed.
  \item[\yamlkey{log\_policy}, \yamlkey{progress\_interval\_sec}] how command output and progress
        are logged.
  \item[\yamlkey{diagnostics\_enabled}, \yamlkey{diagnostics\_interval\_sec}] optional extra
        scheduler diagnostics.
\end{description}

\begin{jwarn}
\yamlkey{per\_task\_timeout\_sec} is a \emph{per-command} limit; a calculator's \yamlkey{timeout}
(Chapter~\ref{ch:calculators}) is a \emph{total} limit for one sample's whole \yamlkey{execution}.
When both are set, the effective per-command limit is bounded by the remaining calculator budget.
\end{jwarn}

\section{Summary}
\label{sec:rt-summary}

\yamlkey{Runtime} is the performance and execution-control block: \yamlkey{mode}/\yamlkey{workers}
for parallelism, \yamlkey{batch\_size} for batching, \yamlkey{writer} for persistence,
\yamlkey{sample\_artifacts} for disk use, and \yamlkey{Subprocess} for external-command scheduling.
Every key defaults to v1 behaviour, so adding a \yamlkey{Runtime} block only ever changes speed and
disk use---never your results. The next chapter puts these to work.
